\documentclass[11pt]{article}
\usepackage[margin=1in]{geometry}
\usepackage[T1]{fontenc}
\usepackage{lmodern}
\usepackage{hyperref}
\usepackage{listings}
\usepackage{longtable}
\usepackage{booktabs}
\usepackage{enumitem}
\usepackage{xcolor}

\hypersetup{colorlinks=true,linkcolor=blue,urlcolor=blue}
\lstset{basicstyle=\ttfamily\small,columns=fullflexible,breaklines=true,
  frame=single,showstringspaces=false,keepspaces=true}
\newcommand{\Current}{\textbf{Current}}
\newcommand{\Proposed}{\textbf{Proposed}}
\newcommand{\Experimental}{\textbf{Experimental}}
\newcommand{\Invariant}{\textbf{Invariant}}

\title{A Python Software-Library Interface for the\\
Chaos Language Algorithm\\
Including Dynamics-Aware Trajectory Embedding}
\author{CLA Project}
\date{July 12, 2026}

\begin{document}
\maketitle

\begin{abstract}
This document specifies a user-facing and extensible Python interface for the
Chaos Language Algorithm (CLA) as both a symbolic grammar learner and an
end-to-end trajectory-analysis pipeline. The pipeline converts a scalar or
vector trajectory into a low-dimensional kinetic representation, partitions it
into a bounded microstate alphabet, and induces an exact-reconstructing grammar
using sequential chunks, substitutable categories, and minimum-description-
length (MDL) search. The design preserves the dependency-free pure-Python
implementation while defining optional NumPy, scikit-learn, deeptime, and
Hyperon adapters. It explicitly distinguishes implemented behavior from proposed
interfaces and from scientifically unvalidated diagnostics. The central API
contract is that fitted symbolic models reconstruct their training symbol stream
exactly; predictive quality, grammar survival, and dynamical interpretation are
separate measured claims, never consequences of reconstruction alone.
\end{abstract}

\tableofcontents

\section{Scope, status, and design goals}

The current repository already implements a deterministic symbolic-string CLA,
fixed-grid and k-means symbolizers, a small pure-Python non-reversible kinetic
map, optional deeptime routines, JSON state persistence and edit replay, and
Phase-0 evaluation diagnostics. These pieces are not yet integrated behind one
fitted trajectory-pipeline object. This specification supplies that integration
contract without pretending it already exists.

Every interface below has one of three status labels.
\begin{itemize}[nosep]
\item \Current: present in the repository, though its exact signature may be
  normalized by a compatibility wrapper.
\item \Experimental: present, but numerically, statistically, or operationally
  provisional.
\item \Proposed: a normative target for future implementation.
\end{itemize}

The design goals are:
\begin{enumerate}[nosep]
\item one simple API for symbolic streams and one for raw trajectories;
\item exact reconstruction of the fitted symbolic stream;
\item explicit separation of chunks (sequences) and categories (alternatives);
\item fitted preprocessing, embedding, and symbolization objects that can
  transform new data without refitting or leaking holdout information;
\item deterministic, seed-controlled operation where the selected backend
  permits it;
\item immutable, serializable domain records and replayable accepted edits;
\item backend-neutral Protocols suitable for pure Python, deeptime, and Hyperon;
\item honest evaluation: baseline diagnostics are labeled as baselines, not as
  a validated CLA likelihood or proof of a strange attractor.
\end{enumerate}

\section{Conceptual pipeline and contracts}

The end-to-end data path is
\begin{verbatim}
raw trajectory
  -> validation and preprocessing fitted on training data
  -> intrinsic-dimension diagnostics
  -> lagged TICA/VAMP kinetic embedding
  -> bounded-k k-means microstates
  -> optional PCCA+ memberships used only as category seeds
  -> symbolic stream
  -> chunk and category proposals
  -> MDL-scored grammar search
  -> fitted model, edit trace, diagnostics, and exact expansion
\end{verbatim}

For high-dimensional data the default is non-reversible VAMP, not reversible
TICA. The lagged covariance must not be silently symmetrized because directional
symbolic structure may be scientifically relevant. PCCA+ macrostates are not the
default base alphabet: microstates remain the observations given to CLA, while
soft metastable memberships may seed category proposals.

Let the training trajectory be $X=(x_0,\ldots,x_{n-1})$. With lag $\tau$, an
embedding may produce only $n-\tau$ coordinates. The fitted artifact therefore
carries an explicit index map from every coordinate and symbol back to a source
time index. No API may imply one symbol per original sample when a transform has
trimmed endpoints.

\paragraph{Exactness boundary.}
If $S_0$ is the symbolic stream supplied to grammar induction and $(G,P)$ is the
fitted grammar and parse, then
\[
  \operatorname{expand}(P,G)=S_0.
\]
This does not reconstruct the floating-point trajectory. A lossy embedding and
partition remain lossy. A proposed quantized trajectory decoder, if added, must
be named separately and report its distortion.

\section{Package and dependency boundaries}

\subsection{Proposed public layout}

\begin{verbatim}
chaoslang/
  __init__.py              # stable exports: CLA, CLAPipeline, configs
  api.py                   # compatibility symbolic facade
  pipeline.py              # end-to-end estimator facade
  model.py                 # fitted model and result bundles
  config.py                # frozen versioned configuration records
  core/                    # dependency-free immutable domain objects
  preprocessing/           # validation, scaling, missing-data policies
  diagnostics/             # intrinsic dimension and data-quality checks
  embedding/               # Protocols and pure-Python kinetic map
  symbolization/           # fitted partition Protocols and pure k-means
  mining/                  # chunk-miner Protocol and implementations
  categorization/          # category-inducer Protocol and implementations
  scoring/                 # code lengths, Scorer Protocol, MDL
  search/                  # proposal scheduling and acceptance
  evaluation/              # surrogates, prediction, sweeps
  persistence/             # versioned artifact schema and replay
  stores/                  # Fact and FactStore seam
  backends/deeptime/       # optional VAMP/TICA/PCCA+ implementation
  backends/hyperon/        # optional fact/search adapters
  benchmarks/              # controls, not core API
\end{verbatim}

For migration, existing single-file modules such as
\texttt{chaoslang.embedding} and \texttt{chaoslang.symbolization} may continue to
re-export implementations. Core must not import NumPy, sklearn, deeptime, or
Hyperon. Optional dependencies are loaded only after an explicitly selected
backend requires them. Importing \texttt{chaoslang} must remain lightweight.

\subsection{Backend capability discovery}

\Proposed:
\begin{lstlisting}[language=Python]
@dataclass(frozen=True)
class BackendInfo:
    name: str
    version: str | None
    capabilities: frozenset[str]
    deterministic_notes: tuple[str, ...] = ()

def available_backends() -> tuple[BackendInfo, ...]: ...
\end{lstlisting}

Selection uses registered names or supplied objects. The value
\texttt{"auto"} may choose a backend only if the resolved backend is recorded in
the fitted provenance. Missing optional software raises
\texttt{BackendUnavailableError} with an installation extra; it must not silently
change algorithms.

\section{Typed public data model}

The public model should use frozen dataclasses and explicit IDs. Existing
\texttt{Token}, \texttt{CategoryOccurrence}, \texttt{Production},
\texttt{Category}, \texttt{Corpus}, \texttt{Grammar}, \texttt{Score},
\texttt{Edit}, and \texttt{GrammarState} remain valid compatibility forms.
The normalized target is:

\begin{lstlisting}[language=Python]
from dataclasses import dataclass, field
from typing import Literal, Mapping, NewType, Sequence, TypeAlias

TokenId = NewType("TokenId", str)
RuleId = NewType("RuleId", str)
CategoryId = NewType("CategoryId", str)
RunId = NewType("RunId", str)
Point: TypeAlias = tuple[float, ...]
Symbol: TypeAlias = str

@dataclass(frozen=True)
class Token:
    id: TokenId
    kind: Literal["base", "chunk"]
    label: str

@dataclass(frozen=True)
class CategoryOccurrence:
    category_id: CategoryId
    member_id: TokenId

ParseEntry: TypeAlias = TokenId | CategoryOccurrence

@dataclass(frozen=True)
class Production:
    id: RuleId
    lhs: TokenId
    rhs: tuple[ParseEntry, ...]

@dataclass(frozen=True)
class Category:
    id: CategoryId
    members: frozenset[TokenId]
    weights: Mapping[TokenId, float] = field(default_factory=dict)

@dataclass(frozen=True)
class SymbolStream:
    symbols: tuple[Symbol, ...]
    source_indices: tuple[int, ...]
    metadata: Mapping[str, object] = field(default_factory=dict)

@dataclass(frozen=True)
class Score:
    model_bits: float
    data_bits: float
    diagnostics: Mapping[str, float] = field(default_factory=dict)
    @property
    def total_bits(self) -> float: ...

@dataclass(frozen=True)
class GrammarState:
    original: SymbolStream
    parse: tuple[ParseEntry, ...]
    grammar: Grammar
    score: Score
    edit_log: tuple[Edit, ...] = ()
\end{lstlisting}

A category is not encoded as a production whose right side lists members. Each
category occurrence stores its observed member so expansion is unambiguous.
IDs, not labels, establish identity. Mapping fields are exposed as read-only
views or persistent maps; freezing a dataclass alone does not make a contained
\texttt{dict} immutable.

\subsection{Trajectory and fitted-stage records}

\Proposed:
\begin{lstlisting}[language=Python]
@dataclass(frozen=True)
class Trajectory:
    values: Sequence[Sequence[float] | float]
    times: Sequence[float] | None = None
    metadata: Mapping[str, object] = field(default_factory=dict)

@dataclass(frozen=True)
class IntrinsicDimensionReport:
    participation_ratio: float
    covariance_eigenvalues: tuple[float, ...]
    variance_dimension: int
    variance_fraction: float
    warnings: tuple[str, ...] = ()

@dataclass(frozen=True)
class EmbeddingResult:
    coordinates: tuple[Point, ...]
    source_indices: tuple[int, ...]
    singular_values: tuple[float, ...]
    timescales: tuple[float, ...] = ()
    metadata: Mapping[str, object] = field(default_factory=dict)

@dataclass(frozen=True)
class PartitionResult:
    stream: SymbolStream
    centers: tuple[Point, ...]
    inertia: float | None
    soft_memberships: tuple[tuple[float, ...], ...] = ()

@dataclass(frozen=True)
class FitTrace:
    scores: tuple[Score, ...]
    candidates_evaluated: tuple[int, ...]
    accepted_edits: tuple[Edit, ...]
    stop_reason: str
\end{lstlisting}

\section{Extension Protocols}

The following Protocols are the stable replaceable seams. Implementations may
use classes, functions wrapped as adapters, local stores, or remote services,
but results must be materialized into the public immutable records.

\begin{lstlisting}[language=Python]
from typing import Iterable, Protocol, Self

class Preprocessor(Protocol):
    def fit(self, x: Trajectory) -> Self: ...
    def transform(self, x: Trajectory) -> Trajectory: ...
    def fit_transform(self, x: Trajectory) -> Trajectory: ...

class DimensionDiagnostic(Protocol):
    def analyze(self, x: Trajectory) -> IntrinsicDimensionReport: ...

class TrajectoryEmbedder(Protocol):
    def fit(self, x: Trajectory) -> Self: ...
    def transform(self, x: Trajectory) -> EmbeddingResult: ...
    def fit_transform(self, x: Trajectory) -> EmbeddingResult: ...

class Symbolizer(Protocol):
    def fit(self, x: EmbeddingResult) -> Self: ...
    def transform(self, x: EmbeddingResult) -> PartitionResult: ...
    def fit_transform(self, x: EmbeddingResult) -> PartitionResult: ...

class PatternMiner(Protocol):
    def proposals(self, state: GrammarState) -> Iterable[Proposal]: ...

class CategoryInducer(Protocol):
    def proposals(self, state: GrammarState,
                  seeds: CategorySeeds | None = None
                  ) -> Iterable[Proposal]: ...

class Scorer(Protocol):
    def score(self, state: GrammarState) -> Score: ...
    def delta(self, state: GrammarState, proposal: Proposal) -> float: ...

class SearchStrategy(Protocol):
    def fit(self, initial: GrammarState, *, rng: RandomSource
            ) -> tuple[GrammarState, FitTrace]: ...

class Predictor(Protocol):
    def predict_proba(self, context: Sequence[Symbol]
                      ) -> Mapping[Symbol, float]: ...

class FactStore(Protocol):
    def add(self, fact: Fact) -> None: ...
    def query(self, predicate: str | None = None) -> Iterable[Fact]: ...
\end{lstlisting}

A miner or inducer returns proposals and never mutates state. Only an edit
applier/search strategy commits edits. The current repository implements a
minimal \texttt{FactStore}, n-gram and bounded suffix-trie miners, exact-context
and Jensen--Shannon category inducers, and a greedy scorer loop; the richer
Protocol signatures above are proposed normalization.

\section{Configuration and deterministic randomness}

Configurations are frozen, JSON-serializable, versioned, and nested. Unknown
fields are errors by default. No mutable estimator object is embedded in a
portable configuration; plugin objects are supplied separately and their
qualified names and versions enter provenance.

\begin{lstlisting}[language=Python]
@dataclass(frozen=True)
class EmbeddingConfig:
    method: Literal["identity", "vamp", "tica"] = "vamp"
    backend: Literal["pure", "deeptime", "auto"] = "pure"
    dimension: int = 5
    lag: int = 1
    shrinkage: float = 1e-6
    reversible: bool = False

@dataclass(frozen=True)
class SymbolizationConfig:
    method: Literal["kmeans", "fixed_grid"] = "kmeans"
    microstates: int = 32
    prefix: str = "km"
    max_iterations: int = 50
    pcca_macrostates: int | None = None

@dataclass(frozen=True)
class GrammarConfig:
    max_iterations: int = 8
    miner: Literal["suffix_trie", "ngram"] = "suffix_trie"
    n_min: int = 2
    n_max: int = 5
    min_uses: int = 2
    categories: Literal["off", "exact", "js"] = "js"
    js_threshold: float = 0.1

@dataclass(frozen=True)
class EvaluationConfig:
    surrogate_runs: int = 8
    heldout_fraction: float = 0.3
    predictor_order: int = 1
    smoothing_alpha: float = 0.5

@dataclass(frozen=True)
class CLAConfig:
    schema_version: str = "1.0"
    seed: int = 0
    embedding: EmbeddingConfig = field(default_factory=EmbeddingConfig)
    symbolization: SymbolizationConfig = field(
        default_factory=SymbolizationConfig)
    grammar: GrammarConfig = field(default_factory=GrammarConfig)
    evaluation: EvaluationConfig = field(default_factory=EvaluationConfig)
\end{lstlisting}

\texttt{reversible=True} is valid only with method \texttt{"tica"}; method
\texttt{"vamp"} is non-reversible. The implementation rejects contradictory
settings rather than guessing.

One root seed is expanded deterministically into named child seeds, for example
\texttt{preprocess}, \texttt{embedding}, \texttt{kmeans}, \texttt{search}, and
\texttt{surrogate/i}. Child generation must use a documented stable hash rather
than Python's process-randomized \texttt{hash()}. Provenance records every child
seed. A backend that cannot guarantee bitwise determinism records that fact and
must not claim it.

\section{Simple API}

\subsection{Symbolic streams}

The existing entry point remains supported:
\begin{lstlisting}[language=Python]
from chaoslang import CLA

model = CLA.simple(max_iterations=8, seed=0).fit_symbols(
    "abcabcabc"
)
assert model.expand() == tuple("abcabcabc")
\end{lstlisting}

\Current: \texttt{CLA.simple(...).fit\_symbols(...)} and
\texttt{CLAModel.expand()}. \Proposed additions:
\begin{lstlisting}[language=Python]
class CLA:
    @classmethod
    def simple(cls, *, seed: int = 0, **kwargs: object) -> "CLA": ...
    def fit(self, symbols: Sequence[str]) -> "SymbolicCLAModel": ...
    def fit_symbols(self, symbols: str | Iterable[str | Token]
                    ) -> "SymbolicCLAModel": ... # compatibility
    def evaluate(self, symbols: Sequence[str], **kwargs: object
                 ) -> EvaluationReport: ...

class SymbolicCLAModel:
    grammar_: Grammar
    state_: GrammarState
    trace_: FitTrace
    def transform(self, symbols: Sequence[str]) -> EncodedParse: ...
    def inverse_transform(self, parse: EncodedParse) -> tuple[str, ...]: ...
    def expand(self) -> tuple[str, ...]: ...
    def predict(self, contexts: Sequence[Sequence[str]]) -> tuple[str, ...]: ...
    def predict_proba(self, contexts: Sequence[Sequence[str]]
                      ) -> tuple[Mapping[str, float], ...]: ...
\end{lstlisting}

\texttt{transform} parses new symbols using the frozen grammar and does not
learn new rules. Unknown-symbol behavior is configured as \texttt{"error"},
\texttt{"literal"}, or a predeclared unknown token. \texttt{inverse\_transform}
expands any valid encoded parse. Prediction is unavailable until a predictor is
fitted; it raises \texttt{NotFittedCapabilityError}, never fabricates a grammar
likelihood.

\subsection{Raw trajectories}

\Proposed:
\begin{lstlisting}[language=Python]
from chaoslang import CLAPipeline, CLAConfig

pipe = CLAPipeline.simple(seed=7, embedding_dim=5,
                          lag=2, microstates=32)
model = pipe.fit(trajectory)
symbols = model.transform(another_trajectory)
report = model.evaluate(another_trajectory)
assert model.symbolic_model_.expand() == model.training_symbols_.symbols
\end{lstlisting}

The convenience estimator has familiar methods:
\begin{lstlisting}[language=Python]
class CLAPipeline:
    def __init__(self, config: CLAConfig, *,
                 preprocessor: Preprocessor | None = None,
                 diagnostic: DimensionDiagnostic | None = None,
                 embedder: TrajectoryEmbedder | None = None,
                 symbolizer: Symbolizer | None = None,
                 miner: PatternMiner | None = None,
                 category_inducer: CategoryInducer | None = None,
                 scorer: Scorer | None = None,
                 search: SearchStrategy | None = None) -> None: ...

    @classmethod
    def simple(cls, *, seed: int = 0, embedding_dim: int = 5,
               lag: int = 1, microstates: int = 32,
               backend: str = "pure") -> "CLAPipeline": ...

    def fit(self, x: TrajectoryLike, y: None = None) -> "CLAPipelineModel": ...
    def transform(self, x: TrajectoryLike) -> SymbolStream: ...
    def fit_transform(self, x: TrajectoryLike, y: None = None
                      ) -> SymbolStream: ...
    def evaluate(self, x: TrajectoryLike | None = None,
                 *, plan: EvaluationPlan | None = None
                 ) -> EvaluationReport: ...
\end{lstlisting}

\texttt{fit\_transform} returns the fitted training symbol stream, not the
compressed parse. The compressed representation is
\texttt{model.symbolic\_model\_.state\_.parse}. This naming prevents an
ambiguous interpretation of ``transform.'' The fitted pipeline model exposes:
\begin{lstlisting}[language=Python]
@dataclass(frozen=True)
class CLAPipelineModel:
    config_: CLAConfig
    preprocessor_: Preprocessor
    diagnostic_: IntrinsicDimensionReport
    embedder_: TrajectoryEmbedder
    symbolizer_: Symbolizer
    training_embedding_: EmbeddingResult
    training_partition_: PartitionResult
    symbolic_model_: SymbolicCLAModel
    provenance_: Provenance

    @property
    def training_symbols_(self) -> SymbolStream: ...
    def transform(self, x: TrajectoryLike) -> SymbolStream: ...
    def encode(self, x: TrajectoryLike) -> EncodedParse: ...
    def predict(self, x: TrajectoryLike) -> tuple[str, ...]: ...
    def predict_proba(self, x: TrajectoryLike
                      ) -> tuple[Mapping[str, float], ...]: ...
    def evaluate(self, x: TrajectoryLike | None = None,
                 *, plan: EvaluationPlan | None = None
                 ) -> EvaluationReport: ...
    def save(self, path: str | Path) -> None: ...
\end{lstlisting}

\section{Advanced end-to-end example}

\begin{lstlisting}[language=Python]
from chaoslang import CLAPipeline
from chaoslang.config import (
    CLAConfig, EmbeddingConfig, SymbolizationConfig,
    GrammarConfig, EvaluationConfig,
)

config = CLAConfig(
    seed=19,
    embedding=EmbeddingConfig(
        method="vamp", backend="deeptime",
        dimension=6, lag=4, reversible=False,
    ),
    symbolization=SymbolizationConfig(
        method="kmeans", microstates=48,
        pcca_macrostates=4, # category seeds, not base symbols
    ),
    grammar=GrammarConfig(
        miner="suffix_trie", n_min=2, n_max=10,
        min_uses=2, categories="js", js_threshold=0.08,
    ),
    evaluation=EvaluationConfig(
        surrogate_runs=20, heldout_fraction=0.3,
        predictor_order=2,
    ),
)

pipeline = CLAPipeline(config)
model = pipeline.fit(train_trajectory)

# Frozen fitted stages transform the test continuation; no refitting.
test_symbols = model.transform(test_trajectory)
encoded = model.encode(test_trajectory)
assert model.symbolic_model_.inverse_transform(encoded) == test_symbols.symbols

report = model.evaluate(test_trajectory)
print(report.surrogate.delta_grammar_bits)
print(report.heldout.log_loss_bits_per_symbol)
print(report.heldout.perplexity)
model.save("lorenz_cla.cla")
\end{lstlisting}

Train/test splitting must occur in trajectory time before fitting scaling,
embedding, clustering, grammar, or predictor. For a single continuous trace,
lagged windows crossing the split are either excluded or explicitly marked as
warm-up context. Random pointwise splitting is invalid for temporal evaluation.

\section{Fitting semantics by stage}

\subsection{Preprocessing}

The default validates finite rectangular data and optionally centers/scales using
training statistics. Missing values are rejected unless an explicit imputer is
configured. Irregular sample times are either rejected by lag-in-steps methods
or resampled by an explicit preprocessor. Input arrays are not mutated.

\subsection{Intrinsic dimension}

\Current/\Experimental: covariance spectral participation ratio and a
variance-fraction suggested dimension. This is a cheap diagnostic, not a
correlation dimension and not proof of an attractor dimension. The report is
advisory unless \texttt{dimension="auto"} is explicitly selected. Automatic
selection records the rule and any clipping bounds.

\subsection{Kinetic embedding}

\Current/\Experimental pure path: shrinkage-regularized $C_{00}$,
$C_{0\tau}$, and $C_{\tau\tau}$; non-symmetrized lagged covariance; singular
values clipped to $[0,1]$; small Jacobi eigensolver. It is suitable for tests and
small controls, not production dimensions near 200--300.

\Current optional path: deeptime VAMP or TICA, singular values, timescales, and
cumulative kinetic variance. The fitted API must preserve and serialize the
actual projection model needed by \texttt{transform}; returning coordinates
alone is insufficient. Lag selection by implied-timescale plateau or
cross-validated VAMP-2 is \Proposed.

\subsection{Microstates and PCCA+}

\Current: deterministic farthest-first pure-Python k-means and optional sklearn
k-means wrappers. The normalized fitted symbolizer stores centers and assigns
new points to nearest centers. It never refits on transform. Labels are stable
IDs based on fitted center order.

\Current/\Experimental optional PCCA+ helper yields microstate assignments and
soft memberships, but its broad exception fallback and lack of a fitted public
object are not sufficient for a stable API. The proposed API surfaces failures
as diagnostics or typed exceptions. PCCA+ memberships become
\texttt{CategorySeeds}; they do not replace microstate symbols by default.

\subsection{Grammar induction}

\Current: bounded n-gram or suffix-trie chunk proposals; exact-context or
Jensen--Shannon category proposals; approximate two-part MDL; deterministic
greedy selection with lexical tie breaking; dead-rule pruning; explicit
\texttt{M[v]} category occurrences; exact reconstruction assertions.

\Proposed search trace records the initial score and every accepted score,
proposal provenance, exact delta, estimated delta, tie-break key, stopping reason,
and runtime counters. Joint chunk/category moves, McBride-derivative pre-ranking,
beam search, and calibrated real-bit coding are future strategies behind the
same Protocols.

\section{Prediction and evaluation}

Evaluation results identify the data split, fitted artifact, metric algorithm,
number of evaluated symbols, seed schedule, and uncertainty samples.

\begin{lstlisting}[language=Python]
@dataclass(frozen=True)
class SurrogateReport:
    real_gain_bits: float
    surrogate_gain_bits: tuple[float, ...]
    surrogate_mean_bits: float
    delta_grammar_bits: float

@dataclass(frozen=True)
class PredictiveReport:
    log_loss_bits_per_symbol: float
    perplexity: float
    evaluated_symbols: int
    vocabulary_size: int
    predictor: str

@dataclass(frozen=True)
class RateDistortionPoint:
    embedding_dimension: int
    microstates: int
    distortion: float | None
    delta_grammar_bits: float
    heldout_log_loss_bits: float

@dataclass(frozen=True)
class EvaluationReport:
    surrogate: SurrogateReport | None
    heldout: PredictiveReport | None
    rate_distortion: tuple[RateDistortionPoint, ...]
    implied_timescales: ImpliedTimescaleReport | None
    warnings: tuple[str, ...]
    provenance: Provenance
\end{lstlisting}

\paragraph{Surrogate excess compression.}
\Current/\Experimental: shuffle symbols while preserving marginal counts, fit
CLA to each surrogate, and subtract mean surrogate gain from real-stream gain.
The current no-grammar baseline is empirical iid code length and the grammar
cost is the current approximate scorer. Therefore values are Phase-0 comparative
diagnostics, not calibrated evidence in final real bits. Surrogate algorithms
must be named; future block, phase-randomized, and dynamical surrogates belong
behind a \texttt{SurrogateGenerator} Protocol.

\paragraph{Held-out loss and perplexity.}
\Current/\Experimental: additive-smoothed n-gram next-symbol predictor. It is
not a CLA grammar likelihood. Reports must say \texttt{predictor="ngram"}. A
future CLA predictor may traverse grammar-derived contexts and must normalize
probabilities over the declared vocabulary. Perplexity is $2^{\ell}$ for loss
$\ell$ in bits per symbol.

\paragraph{Rate--distortion.}
\Proposed: refit complete pipelines over embedding dimension $d$ and microstate
count $k$, reporting predictive loss, surrogate excess compression, and an
explicit distortion measure. Every grid point has an independent run record;
shared preprocessing is allowed only when it cannot leak validation data.

\paragraph{Implied timescales.}
\Current optional helper computes timescales across lags. Plateau selection and
confidence intervals are proposed. Timescales can diagnose lag choice but do not
by themselves validate CLA grammar.

\section{Persistence, replay, and provenance}

\Current persistence round-trips \texttt{GrammarState} as deterministic JSON,
renders stable text, and replays the current edit log for covered edits.
\Proposed artifact persistence extends this to all fitted stages.

A saved \texttt{.cla} artifact is a directory or ZIP with:
\begin{verbatim}
manifest.json             schema, hashes, package/backend versions
config.json               resolved immutable configuration
provenance.json           run ID, seeds, platform, warnings
preprocessor.json         fitted statistics
embedder/                 portable parameters or backend payload
symbolizer.json           centers, labels, optional memberships
grammar-state.json        corpus identity, parse, grammar, scores, edits
fit-trace.json             all accepted transitions and stop reason
evaluation/*.json         reports with data-split identities
\end{verbatim}

\begin{lstlisting}[language=Python]
@dataclass(frozen=True)
class Provenance:
    run_id: RunId
    created_utc: str
    chaoslang_version: str
    artifact_schema_version: str
    config_sha256: str
    input_sha256: str | None
    backend_versions: Mapping[str, str]
    root_seed: int
    child_seeds: Mapping[str, int]
    platform: Mapping[str, str]
    warnings: tuple[str, ...] = ()

class CLAPipelineModel:
    def save(self, path: str | Path, *, include_training_data: bool = False
             ) -> None: ...
    @classmethod
    def load(cls, path: str | Path, *, verify_hashes: bool = True,
             allow_unsafe_backend_payload: bool = False
             ) -> "CLAPipelineModel": ...
\end{lstlisting}

JSON and numeric arrays are preferred over pickle. Unsafe backend payloads are
disabled by default. Loading validates schema version, hashes, dimensions,
finite values, grammar references, category membership, edit sequence, and exact
symbolic reconstruction when the original symbols are included. Replay applies
serialized typed edits at recorded state hashes; it does not search for a
similar block and silently skip unmatched edits. The current best-effort replay
behavior remains a legacy reader only.

Training data are excluded by default for size and privacy. If excluded, the
artifact stores a content hash, shape, source-index range, and metadata sufficient
to identify the data externally. Migration functions are explicit and never
rewrite the source artifact in place.

\section{Errors, validation, and invariants}

\Proposed public exception hierarchy:
\begin{verbatim}
CLAError
  ConfigurationError
  InvalidTrajectoryError
  InsufficientDataError
  NotFittedError
  NotFittedCapabilityError
  BackendUnavailableError
  BackendComputationError
  UnknownSymbolError
  InvariantViolation
  ArtifactError
    ArtifactIntegrityError
    UnsupportedSchemaError
    ReplayMismatchError
\end{verbatim}

Required checks include:
\begin{enumerate}[nosep]
\item trajectory points have one positive, constant dimension and finite values;
\item sample count exceeds lag and any estimator-specific minimum;
\item $1\leq d\leq D$ where applicable and $1\leq k\leq n_{embedded}$;
\item source indices are strictly increasing and align with coordinates/symbols;
\item category member sets contain at least two distinct known token IDs;
\item production and parse references resolve; chunk productions are acyclic;
\item proposed occurrence ranges are valid and non-overlapping;
\item accepted edits do not increase MDL under a strict greedy policy;
\item after every accepted edit and after load/replay, expansion equals the
  original symbolic stream;
\item fitted \texttt{transform} never updates means, projections, centers,
  grammar, or random state;
\item probability predictions are finite, non-negative, and sum to one within
  tolerance.
\end{enumerate}

Assertions may aid development but public input and artifact validation must use
typed exceptions because Python can disable assertions.

\section{Hyperon and fact-layer seam}

\Current: a minimal backend-neutral \texttt{Fact} and \texttt{FactStore}, plus
memory storage and state/fact tests. \Proposed: a lossless versioned projection
from domain objects and proposals into facts.

\begin{lstlisting}[language=Python]
@dataclass(frozen=True)
class Fact:
    predicate: str
    args: tuple[object, ...]
    truth: float = 1.0
    attrs: Mapping[str, object] = field(default_factory=dict)

class StateProjector(Protocol):
    def to_facts(self, state: GrammarState) -> Iterable[Fact]: ...
    def from_facts(self, facts: Iterable[Fact]) -> GrammarState: ...
\end{lstlisting}

Hyperon miners and category inducers query a fact store and return ordinary
\texttt{Proposal} records. They cannot mutate the authoritative grammar state.
An ordinary edit applier revalidates all references, computes exact score, and
checks reconstruction before commit. Hyperon truth values and provenance may
rank proposals but do not bypass MDL acceptance. This keeps the scientific
result portable and testable without a Hyperon runtime.

\section{Backward compatibility and migration}

The compatibility policy for the next major implementation stage is:
\begin{enumerate}[nosep]
\item retain \texttt{from chaoslang import CLA};
\item retain \texttt{CLA.simple(...).fit\_symbols(...)};
\item retain \texttt{CLAModel.state}, \texttt{score}, \texttt{grammar},
  \texttt{parse}, \texttt{expand\_tokens()}, and \texttt{expand()};
\item re-export current functions from their old module paths;
\item introduce \texttt{CLAPipeline} rather than changing \texttt{CLA.fit} to
  guess whether nested sequences are tokens or trajectory vectors;
\item add deprecation warnings for constructor strings only after equivalent
  config objects and migration documentation are released;
\item load existing grammar-state JSON through a legacy schema reader.
\end{enumerate}

There are current inconsistencies to normalize carefully: compatibility modules
and package modules expose overlapping miner/scorer/store names; the current
\texttt{CLAModel} does not retain seed/config; pure kinetic mapping returns only
coordinates and summary statistics rather than a reusable fitted transform;
optional deeptime routines are functions rather than estimators; and current
replay may search or skip when recorded edits do not match. None should be
silently presented as satisfying the proposed fitted-pipeline contract.

\section{Implementation roadmap}

\subsection*{Stage 1: unify stable symbolic API}
\begin{itemize}[nosep]
\item Freeze and export Protocols, configs, exception types, and versioned IDs.
\item Add fit trace and store resolved seed/config on symbolic models.
\item Preserve all current imports and exact-reconstruction tests.
\item Replace assertion-only public validation with typed exceptions.
\end{itemize}

\subsection*{Stage 2: fitted trajectory stages}
\begin{itemize}[nosep]
\item Implement fitted preprocessor, pure kinetic-map estimator, and fitted
  k-means symbolizer with source-index alignment.
\item Build \texttt{CLAPipeline} and \texttt{CLAPipelineModel}.
\item Add train/holdout leakage tests and transform-without-refit tests.
\end{itemize}

\subsection*{Stage 3: robust optional backend}
\begin{itemize}[nosep]
\item Wrap deeptime VAMP/TICA as fitted embedders with serializable parameters.
\item Convert PCCA+ output into explicit category seeds and remove broad silent
  exception handling.
\item Add implied-timescale and VAMP-2 lag/dimension diagnostics.
\end{itemize}

\subsection*{Stage 4: artifact and evaluation contract}
\begin{itemize}[nosep]
\item Implement portable artifact manifest, hashes, migration, strict replay,
  and provenance capture.
\item Separate n-gram baseline prediction from a future CLA predictor.
\item Add reproducible $d,k$ rate--distortion sweeps and multi-seed intervals.
\end{itemize}

\subsection*{Stage 5: algorithmic extensions}
\begin{itemize}[nosep]
\item Calibrate all description lengths to real bits.
\item Add McBride-derivative pre-ranking, joint chunk/category moves, and beam
  search behind existing Protocols.
\item Add Hyperon adapters only after lossless fact projection and parity tests.
\item Consider probabilistic symbols and learned embeddings only as separate,
  versioned CLA variants.
\end{itemize}

\section{Testing and scientific acceptance}

Unit and property tests should cover type validation, deterministic named seeds,
chunk/category identity, acyclic expansion, replay, source-index alignment,
fitted-stage immutability, and optional-backend parity on small controls.
Integration tests should run
\begin{verbatim}
trajectory -> VAMP/TICA -> k-means -> CLA -> expansion -> persistence -> load
\end{verbatim}
and require the same symbols, grammar, parse, scores (within documented numeric
tolerance), and provenance identities after loading.

Scientific acceptance is stronger than software acceptance. The initial battery
should include periodic and non-chaotic controls, shuffled and structure-aware
surrogates, Lorenz-63 and Roessler, and high-dimensional noisy lifts. It should
report multiple seeds, uncertainty, raw fixed-grid failure, adaptive-partition
results, held-out baselines, and grammar rate--distortion curves. A passing unit
test establishes software behavior only. Positive surrogate excess compression
under the current approximate scorer is suggestive, not conclusive. Exact
reconstruction is a necessary representation invariant, not evidence that a
learned grammar is causal, unique, predictive, or scientifically meaningful.

\section{Current evidence and unresolved choices}

This specification was grounded in the repository's public API, immutable core
records, edit applier/search loop, n-gram and suffix-trie miners, exact and
Jensen--Shannon category induction, symbolization, pure and deeptime embedding
modules, evaluation harness, persistence/replay code, README, project decisions,
high-dimensional embedding design, and tests. At inspection time the branch was
\texttt{agent/hd-embedding-cla} at commit \texttt{16fe6bc}; the full unittest
suite was also invoked as a validation gate for this design work.

The principal unresolved design choices are:
\begin{enumerate}[nosep]
\item the portable serialization form for fitted deeptime projections;
\item the definitive real-bit CLA code and predictive distribution;
\item treatment of multiple independent trajectories versus one concatenated
  stream (boundaries must never create artificial chunks);
\item online/streaming covariance, partition updates, and grammar updates;
\item lag/dimension model selection and confidence-interval methodology;
\item whether category weights remain proposal metadata or become part of a
  probabilistic grammar variant;
\item exact semantics for parsing unseen streams when several rewrites overlap.
\end{enumerate}

These ambiguities should be resolved through versioned interfaces and tests,
not hidden behind convenience defaults.

\section{Summary}

The recommended library has two explicit front doors: \texttt{CLA} for existing
symbolic streams and \texttt{CLAPipeline} for trajectories. The latter owns the
entire fitted path from validation and dynamics-aware embedding through bounded
microstates and grammar induction. Every stage has a typed replaceable Protocol,
every stochastic action derives from a recorded root seed, every accepted edit
is replayable, and every fitted grammar reconstructs its training symbol stream.
Optional deeptime and Hyperon components remain adapters. Most importantly, the
API separates representation exactness from predictive and scientific claims,
so future CLA research can improve algorithms without weakening the meaning of
reported results.

\end{document}
